\section{INTRODUCTION}
\label{sec:introduction}
The practical appeal of O-ISAC lies in reuse. Optical networks already deploy
sources, propagation paths, receivers, and signal processors, and these
resources can reveal information about the environment while carrying data.
Distributed sensing over fiber, positioning based on visible light, and
turbulence monitoring illustrate the breadth of this opportunity
\cite{OISAC_SCR00013,OISAC_SCR00196,OISAC_SCR01085}.

Sharing turns that opportunity into a systems question. A resource or receiver
chosen for one function can change the information available to the other.
The central engineering problem is what can be shared and how that sharing
shapes system design.

This question has gained importance since integrated sensing and communication
was identified as one of the usage scenarios for International Mobile
Telecommunications for 2030 and beyond (IMT-2030)
\cite{ITUR_M2160_2023}. Optical implementations span guided fiber, free space
optical and visible light communication links, systems in the millimeter wave
and terahertz bands enabled by photonics, and hybrid designs
\cite{OISAC_SCR00916,OISAC_SCR00036,OISAC_SCR00083}. Their propagation paths,
observables, hardware, and performance measures differ. O-ISAC therefore forms
a family of related architectures, each with its own physical and measurement
context.

That diversity requires a clear survey boundary and careful comparison. The
survey includes optical or photonic systems that consider, integrate, design,
optimize, or evaluate communication and sensing together. Their common setting
may lie in an architecture, optical link, signal or resource framework,
hardware platform, channel model, or application scenario. Coupling can appear
through a shared waveform, concurrent operation, or another common system
choice. Each result is interpreted through its metric definition, measurement
plane, and operating conditions. This context supports consistent
interpretation and bounded comparison across platforms.

\subsection{Prior Syntheses and Survey Positioning}
\label{sec:related_surveys}

Prior surveys organize complementary parts of O-ISAC. Each serves a different
reader task and uses a different evidence unit and comparison boundary.
Table~\ref{tab:prior_surveys} provides a compact map of these perspectives.

\begin{table*}[!t]
\caption{Representative prior O-ISAC syntheses organized by reader task,
evidence unit, and comparison boundary}
\label{tab:prior_surveys}
\centering
\begingroup
\footnotesize
\setlength{\tabcolsep}{2.2pt}
\renewcommand{\arraystretch}{1.00}
\begin{tabularx}{\textwidth}{@{}%
>{\raggedright\arraybackslash}p{0.112\textwidth}
>{\raggedright\arraybackslash}p{0.140\textwidth}
>{\raggedright\arraybackslash}p{0.120\textwidth}
>{\raggedright\arraybackslash}p{0.190\textwidth}
>{\raggedright\arraybackslash}p{0.160\textwidth}
>{\raggedright\arraybackslash}X@{}}
\toprule
Synthesis family & Sources & Optical scope & Evidence unit and comparison logic &
Reader task & Boundary \\
\midrule
Architecture and system design &
\cite{Wen2024OISACArchitectures,Su2025OISACTheoryPractice,
Chu2025ROISACOverview,Wang2026CrossLayerISACoF} &
General O-ISAC, free space optical links, retroreflective links, and fiber
networks &
Architecture, optimization, target, or link model, with taxonomy and relations
within each scenario &
How are communication and sensing coupled? &
Relations remain tied to each architecture and its assumptions. \\
\addlinespace[1pt]
Analytical limits and estimation methods &
\cite{Khorasgani2026GeneralOpticalISACSurvey,Rode2025PolarizationISAC} &
General optical systems and polarization sensing in coherent fiber links &
Model, bound, objective, estimator, or learning method, with comparison within a
model and among methods &
What limits follow from the assumptions? &
Empirical interpretation remains tied to the model, modality, and evidence
setting. \\
\addlinespace[1pt]
Technological evolution, standardization, and deployment &
\cite{Zhang2026ISACEvolution,Liu2025OpticalNetworkISACProspect,
Cao2026SubmarineFiberISAC,Zhang2026OpticalTransmissionNetworks,
Ip2026DeployedTelecomCableISAC} &
Evolution from radio to optical systems across terrestrial and submarine
networks &
Historical phase, standardization path, network architecture, or deployment,
synthesized by period and ecosystem &
How has O-ISAC moved toward deployment? &
Infrastructure comparison requires aligned physical and operational
conditions. \\
\addlinespace[1pt]
Physical media and optical networks &
\cite{Liang2024LiSACReview,Zhang2023OpticalNetworkISACFrontier,
He2025FiberISACAdvances,Liu2026OpticalFiberISACReview,
Lu2026ConvergedFiberISACReview,Wang2026FiberNetworkISAC} &
Visible light integrated sensing and communication, guided fiber systems, and
optical networks &
Medium, observable, sensing path, or network configuration, with taxonomy within
related solution classes &
How does each path support sensing? &
Transfer requires aligned media, observables, paths, and conditions. \\
\addlinespace[1pt]
Waveforms, resource multiplexing, and photonic hardware &
\cite{Yang2025DSCMISACReview,He2026PhotonicFiberISAC,
Yu2023PhotonAssistedTHz,Bai2025MicrowavePhotonicsISAC,
Wang2025MicrowavePhotonicsISAC,Lyu2026PhotonicTHzWaveforms} &
Fiber, microwave photonics, and millimeter wave or terahertz systems enabled by
photonics &
Waveform, multiplexing, front end, or processing chain, with a mechanism
taxonomy and selected demonstrations &
Which mechanisms enable joint operation? &
Comparison across photonic systems follows alignment of paths, observables, and
metrics. \\
\addlinespace[1pt]
Applications and challenges across platforms &
\cite{Mohsan2026OISACReview} &
Visible light communication, free space optical, fiber, photonics, and hybrid
settings &
Application, use case, solution class, or challenge, with thematic synthesis
across platforms &
Which applications and gaps recur? &
Interpretation across platforms follows the measurement and operating context
of each result. \\
\bottomrule
\end{tabularx}
\endgroup
\end{table*}

Across these perspectives, synthesis across platforms depends on retaining the
physical, measurement, and validation setting of each primary result. This
survey addresses that task through 206 unique O-ISAC studies first made publicly
available from 1 January 2020 through 22 June 2026.

The survey makes four contributions.

\begin{itemize}
\item A taxonomy separates physical modality from the mechanisms that couple
communication and sensing.

\item A comparison framework connects each performance result to its
measurement definition and operating conditions.

\item A synthesis of recurring tradeoff mechanisms retains the conditions that
shape communication and sensing outcomes.

\item An assessment of validation, benchmark readiness, and reproducibility
grounds research priorities in the reviewed evidence.
\end{itemize}

The organization of the article follows the analytical sequence of the survey.
Sections~\ref{sec:foundations_comparison} and~\ref{sec:review_methods}
establish the comparison framework across optical platforms and describe the
review methodology and evidence base. Sections~\ref{sec:optical_modality_families}
through~\ref{sec:technologies_applications_6g} present the technical synthesis
of optical platforms, integration architectures, performance metrics, joint
design tradeoffs, validation, enabling technologies, and applications.
Section~\ref{sec:discussion_roadmap} discusses the research priorities and
limitations that emerge from this synthesis, and
Section~\ref{sec:conclusion} concludes the article.
